\subsection{自伴算子的初等判别法}

命题 9. --- 设 \(v \in \mathcal{L}(\mathrm{F};\mathrm{E})\) 是从 F 到 E 中的连续单射线性映射，其像在 E 中稠密。 \(v\) 的伴随是从 E 到 F 中的连续单射线性映射，且有 \((v^{*})^{-1} = (v^{-1})^{*}\) 。特别地，若 E = F ，则自同态 \(v\) 是埃尔米特的，当且仅当部分算子 \(v^{-1}\) 是自伴的。

部分算子 \(v^{-1}\) 是从 E 到 F 中的闭稠定算子（IV, p. 228 的例 2），而 v 的伴随 \(v^{*}\) 是从 E 到 F 中的连续线性映射；它是单射的，因为 v 的像在 E 中稠密（EVT, V, p. 41，命题 4）。设 s （相应地， \(s'\) ）是从 E ⊕ F 到 F ⊕ E 上的等距同构 \((x,y)\mapsto(-y,x)\) （相应地，从 F ⊕ E 到 E ⊕ F 上的等距同构 \((y,x)\mapsto(-x,y)\) ），并设 \(\iota\) （相应地， \(\iota'\) ）是从 F ⊕ E 到 E ⊕ F 上的等距同构 \((y,x)\mapsto(x,y)\) （相应地，从 E ⊕ F 到 F ⊕ E 上的等距同构 \((x,y)\mapsto(y,x)\) ）。这时有 \(s\circ\iota=-\iota'\circ s'\) ，由此

\[
\Gamma_ {(v ^ {- 1}) ^ {*}} = s \left(\Gamma_ {v ^ {- 1}}\right) ^ {\circ} = s \left(\iota \left(\Gamma_ {v}\right)\right) ^ {\circ} = - \iota^ {\prime} \left(s ^ {\prime} \left(\Gamma_ {v}\right)\right) ^ {\circ} = - \iota^ {\prime} \left(\Gamma_ {v ^ {*}}\right) = \Gamma_ {(v ^ {*}) ^ {- 1}}
\]

（依据 IV, p. 236 的命题 7）。命题得证。

命题 10（海林格–特普利茨）. --- 设 u 是希尔伯特空间 E 上的对称部分算子。若 u 的定义域等于 E ，则 \(u \in \mathcal{L}(E)\) ，且 u 是埃尔米特的。

事实上，部分算子 u 是可闭的（IV, p. 237 的命题 8），而由于它的定义域是 E ，它是闭的（IV, p. 228，注）。于是引用 EVT, I, p. 19，推论 5 即得结论。

推论. --- 设 u 是 E 上的对称部分算子。若部分算子 u 诱导一个从 dom(u) 到 E 上的双射线性映射，则 u 是自伴的。

事实上， u 的逆部分算子 \(u^{-1}\) 是对称的且以 E 为定义域（命题 9），于是 \(u^{-1}\) 是 \(\mathcal{L}(\mathrm{E})\) 中的自伴元素（命题 10），从而 u 是埃尔米特的（命题 9）。

命题 11. --- 设 u 是 E 上的对称部分算子，并设 \(\lambda \in \mathbf{C}\) 。若 \(u + \lambda 1_{\mathrm{E}}\) 与 \(u + \overline{\lambda} 1_{\mathrm{E}}\) 都是满射的，则 u 是自伴的。

只须证明 \(\operatorname{dom}(u^{*}) \subset \operatorname{dom}(u)\) 。设 \(x \in \operatorname{dom}(u^{*})\) 。由假设，存在 \(y \in \operatorname{dom}(u)\) 使得 \(u(y) + \overline{\lambda}y = u^{*}(x) + \overline{\lambda}x\) 。我们来证明 y = x 。对所有 \(z \in \operatorname{dom}(u)\) ，得到

\[
\begin{array}{r l} & {\langle (u + \lambda 1 _ {\mathrm{E}}) (z) | x \rangle = \langle z | (u ^ {*} + \overline {{\lambda}} 1 _ {\mathrm{E}}) (x) \rangle} \\ & {\qquad = \langle z | (u + \overline {{\lambda}} 1 _ {\mathrm{E}}) (y) \rangle = \langle (u + \lambda 1 _ {\mathrm{E}}) (z) | y \rangle ,} \end{array}
\]

（因为 u 是对称的）。由于算子 \(u + \lambda1_{E}\) 是满射的，确有 y = x ，从而 \(x \in \text{dom}(u)\) 。

我们稍后将看到（参见 IV, p. 248 的命题 17）：若 \(\lambda \in \mathbf{C} - \mathbf{R}\) ，则逆命题成立。

命题 12. --- 设 u 是从 E 到 F 中的闭稠定算子。 E 上的部分算子 \(u^{*}\circ u\) 是自伴且正的。它的定义域是 u 的一个核。

用 \(v\) 表示部分算子 \(1_{\mathrm{E}} + u^{*} \circ u\) 。它的定义域是 \(\operatorname{dom}(u^{*} \circ u)\) ，包含在 \(\operatorname{dom}(u)\) 中。对所有 \(x \in \operatorname{dom}(u)\) 与 \(y \in \operatorname{dom}(v)\) ，我们有 \(u(y) \in \operatorname{dom}(u^{*})\) 以及

\begin{enumerate}
\def\labelenumi{(\arabic{enumi})}
\setcounter{enumi}{2}
\tightlist
\item
\end{enumerate}

\[
\langle x \mid v (y) \rangle = \langle x \mid y \rangle + \langle x \mid (u ^ {*} \circ u) (y) \rangle = \langle x \mid y \rangle + \langle u (x) \mid u (y) \rangle ,\tag{4}
\]

\[
\langle y \mid v (y) \rangle = \| y \| ^ {2} + \| u (y) \| ^ {2}.
\]

公式 (4) 蕴含部分算子 \(v\) 是单射的。此外，由 IV, p. 236 的命题 7 \(a\) )，有 \(\mathrm{F} \oplus \mathrm{E} = \Gamma_{u^*} \oplus s(\Gamma_u)\) 。设 \(x \in \mathrm{E}\) 。存在 \(y' \in \mathrm{dom}(u^*)\) 与 \(x' \in \mathrm{dom}(u)\) 使得

\[
(0, x) = (y ^ {\prime}, u ^ {*} (y ^ {\prime})) + (- u (x ^ {\prime}), x ^ {\prime}) = (y ^ {\prime} - u (x ^ {\prime}), x ^ {\prime} + u ^ {*} (y ^ {\prime})).
\]

因此我们有 \(y' = u(x')\) ，从而 \(x' \in \text{dom}(u^* \circ u) = \text{dom}(v)\) 以及

\[
x = x ^ {\prime} + u ^ {*} (y ^ {\prime}) = x ^ {\prime} + (u ^ {*} \circ u) (x ^ {\prime}) = v (x ^ {\prime}).
\]

因此， E 上的部分算子 v 是满射的，并诱导一个从 dom(v) 到 E 上的双射线性映射。

设 \(x \in \mathrm{E}\) 与 \(\operatorname{dom}(v)\) 正交。把 x 写成 \(x = v(x')\) ，其中 \(x' \in \operatorname{dom}(v)\) 。由公式 (4) ，我们得到

\[
0 = \langle x ^ {\prime} | x \rangle = \langle x ^ {\prime} | v (x ^ {\prime}) \rangle = \| x ^ {\prime} \| ^ {2} + \| u (x ^ {\prime}) \| ^ {2}
\]

由此 \(x' = 0\) ，进而 \(x = 0\) 。因此 \(v\) 的定义域在 E 中稠密。

部分算子 \(v\) 稠定；它是双射的，且公式 (3) 表明它是对称的。对 \(v\) 应用命题 10 的推论，我们断定它是自伴的。于是 \(u^{*} \circ u = v - 1_{\mathrm{E}}\) 是自伴的。此外，公式

\[
\langle x \mid (u ^ {*} \circ u) (x) \rangle = \| u (x) \| ^ {2}
\]

（对所有 \(x \in \text{dom}(u^{*} \circ u)\) ）蕴含 \(u^{*} \circ u\) 是正的。

最后，设 \(y \in \mathrm{E}_u\) 与 \(\operatorname{dom}(u^* \circ u)\) 正交。存在 \(x\) 的定义域中的元素 \(u^* \circ u\) 使得 \(y = v(x) = x + (u^* \circ u)(x)\) 。这时我们有 \(0 = (x \mid y)_u = \langle y \mid y \rangle\) ，由此 \(y = 0\) 。

